\scanpage{21}

% L1
\deleted{\uline{\unsure{Nécessité}} de \unsure{faire}}

% L2–9
\uline{Re NB} Si on veut seulement montrer que \deleted{les} topos $R$ est ``algébrique'' c'est encore plus court : On choisit $C$ petite avec \uline{lims} finies, telle que $R\simeq\widetilde C$, \deleted{et} on \deleted{écrit que} part de la structure algébrique $T_C^{\Delta{}_0}$ -- sa restriction \unsure{,} $\Delta$ est donc algébrique, elle est représentée par $\widehat C$, cf \deleted{\unsure{P}} n° 1 -- puis on rajoute des axiomes exprimant la \uline{contin}uité d'un foncteur $C-E$ (commutant déjà aux \uline{lims} finies) par le fait de transformer familles couvrantes en \unsure{famille|familles} couvrantes ...

% L10
\deleted{Je \unsure{note à présent en résumé} que l'algébricité}

% L11
\deleted{\uline{Nécessité}} Il faut suivre pas à pas la construction transfinie

% L12–14
a) Structure vide sur ens de base \inserted{$I$}. Représentée par $\widehat{\widetilde I}$ (cf Pb 1), \inserted{où $\widetilde I=\rho_{\Delta_0}(I)$.} \deleted{$I$ engendre $\Delta$-engendre} \uline{NB} $I$ $\Delta$-engendre $\widehat{\widetilde I}$ en deux pas : rajouter des produits, puis des \uline{lim} quelc.

% L15–17
b) \deleted{\unsure{Ra}} Mettre axiomes sur théorie algébrique $T_R$, $R$ un topos, cela relève du Pb 2 traité. On trouve $R'$ et $R-R'$ \unsure{ess.|en} surjectif (et aussi surjectif sur les Hom).

% L18–22
c) Rajouter flèches (petit ensemble \inserted{$H$}) : \deleted{\unsure{Re.|$R$}} théorie algébrique $T_R$, $R$ \deleted{déjà} topos. \deleted{Soit $C$ l'ens des} Soit \deleted{$C\subset R$ la petite sous-catég} $\mathbf I$ \inserted{$\subset\operatorname{Ob}R$} l'ens des sources et buts de ces flèches, $C$ la catégorie libre engendrée par \deleted{les ob} $I$ comme objets et les flèches $H$, $\rho_\Delta(C)$ la catégorie enveloppante, alors la théorie $T'$ est le produit 2-fibré
% L23
\[
 T'\simeq T_R\mathop{\times}^{2}_{\tau_{\mathbf I}}\tau_C
 \qquad \deleted{\unsure{\tau_I=\underline{\mathrm{Hom}}_{\mathrm{Cat}}}}
\]
% L24–26
où $\tau_C(E)\simeq\underline{\mathrm{Hom}}_{\mathrm{Cat}}(C,E)$, et $\tau_{\mathbf J}(E)$ se définit de $\overline{\unsure{\mathrm m}}$, $J$ considéré comme catégorie \unsure{discrèt.|discrète}. Or (pb 1) $\tau_C\simeq T_{R_1}$, $\tau_I\simeq T_S$ $R_1$ et $S$ deux topos, donc il faut prouver l'existence
% archive sheet number: 20
